\section{The actual rank-four tower and automorphic endpoint}
\label{sec:sr-rank-four-actual-endpoint}

\subsection{The character witnesses on the principal Case II tower}
\label{subsec:sr-case-two-tower-witnesses}
We align the three tower conditions used above with their
actual principal-packet producers.
Here \(\ell\) is odd, \(\ell\nmid N\), and the nontrivial
principal character \(\psi\) has \(2\)-power order and
\(\psi(-1)=(-1)^a\).
Write \(\omega\) for the Teichm\"uller cyclotomic character,
\(\xi\) for the chosen primitive odd auxiliary character
of conductor \(2^n\), and
\[
            \bar\Psi=\bar\omega^{\,a}\bar\psi .
 \tag{T1}
\]
The actual residual diagonal, before the possible
ordinary permutation, is
\[
                  (\bar\xi\bar\Psi,\bar\Psi,\bar\xi,1).
 \tag{T2}
\]
Reduction is injective on all \(2\)-power character values.

Choose an exponent \(s\) with \(\psi^{2^s}=1\), as in
the source's \(\mathrm{psi\_exp}\), and let \(2^c\)
be the conductor of its dyadic component.
The actual choice \(n=\mathrm{nII}(\ell,\psi)\)
satisfies the two numerical lower bounds
\[
       2^n\ge\ell^2 2^{s+2},
                   \qquad 2^n\ge2^{c+6}.
 \tag{T3}
\]
These are precisely the two inequalities extracted from
that definition in
\srcid{ExistsDecompositionElementAtLXiPowFourNeOne.lean}
and
\srcid{ExistsInertiaElementAtTwoXiPowFourNeOne.lean}.
In \(\mathrm{psi\_exp}\) the annihilator is written
\(2^s u\) with \(u\) odd; since our packet has \(2\)-power
order, \(\psi^{2^s}=1\) follows.
No \(\ell\)-primary part of \(\psi\) needs to be killed
in this specialization.

\begin{srlemma}\label{lem:sr-primitive-dyadic-frobenius}
Let \(\xi\) be primitive of conductor \(2^n\), \(n\ge3\).
For an odd integer \(z>1\), if
\(\xi(z)^{2^b u}=1\) with \(u\) odd, then
\[
                              2^n<z^2 2^b .
 \tag{T4}
\]
\end{srlemma}
\begin{proof}
The character is faithful on
\((1+4\mathbb Z_2)/(1+2^n\mathbb Z_2)\).
Its value has \(2\)-power order, so the odd exponent
\(u\) can be removed.
For \(b=0\), squaring gives
\(z^2\equiv1\pmod{2^n}\), hence
\(2^n\le z^2-1<z^2\).
For \(b\ge1\), faithfulness gives
\(z^{2^b}\equiv1\pmod{2^n}\).
The elementary lifting-the-exponent identity yields
\[
 v_2(z^{2^b}-1)=v_2(z^2-1)+b-1.
\]
Consequently
\(2^n\le2^{b-1}(z^2-1)<2^b z^2\).
\end{proof}

At \(\ell\), choose the local reciprocity lift
\(\sigma_\ell\) of a uniformizer with
\(\epsilon_\ell(\sigma_\ell)=1\).
The two finite characters are unramified there, and
their values are \(\psi(\ell),\xi(\ell)\).
Then
\[
       g_\ell=\sigma_\ell^{2^s}\in D_\ell,\qquad
       \omega(g_\ell)=1,\qquad\psi(g_\ell)=1,\qquad
                              \bar\xi(g_\ell)^4\ne1 .
 \tag{T5}
\]
The last assertion follows from (T3)--(T4), with \(z=\ell\)
and \(b=s+2\).
At \(2\), choose an inertia element whose cyclotomic
projection modulo \(2^n\) is \(3\).
Raise it to \(2^c\), obtaining an element \(u\) with
\[
       u\in I_2,\quad \omega(u)=1,\quad\psi(u)=1,\quad
                                   \bar\xi(u)^4\ne1 .
 \tag{T6}
\]
The dyadic part of \(\psi\) is killed, and its other
components are unramified at \(2\).
If the final assertion failed, (T4) with \(z=3\) would
give \(2^n<9\cdot2^{c+2}\), contradicting
\(2^n\ge16\cdot2^{c+2}\).
These arguments also work after an odd
\(\ell\)-power is inserted in the exponents.
They are the actual local witnesses used by the source.

\subsubsection{The finite abelian base field}
Choose three distinct fresh odd primes, outside the
original ramification set \(S\), and quadratic characters
\(\phi,\eta,\eta'\), respectively ramified only at those
primes, with
\[
 \phi(-1)=1,\qquad \eta(-1)=\eta'(-1)=-1 ,
\]
all split at \(S\), and with \(\eta,\eta'\) also split
at the earlier chosen fresh primes.
Such characters are obtained by choosing their prime
conductors in arithmetic progressions.
For the even character choose the conductor prime
\(1\bmod8\); for an odd one choose it \(7\bmod8\).
At each prescribed odd prime require the signed conductor
to be a square modulo that prime.
The Chinese remainder theorem gives a reduced progression
and Dirichlet's theorem supplies a fresh prime in it.
The quadratic characters then have exactly the asserted
ramification and splitting.

Put
\[
 H_0=\{g:\bar\Psi(g)=\phi(g),\
                   \eta(g)=\eta'(g)\},\quad
 F_0=\overline{\mathbb Q}^{\,H_0},\quad
 K=\overline{\mathbb Q}^{\,\ker\eta},\quad L_0=F_0K .
 \tag{T7}
\]
All characters in (T7) are finite and continuous.
Their common kernel contains the commutator, so \(F_0\)
is finite abelian.
Complex conjugation belongs to \(H_0\), by
\(\bar\Psi(c)=1\), \(\phi(c)=1\), and
\(\eta(c)=\eta'(c)=-1\).
Thus \(F_0\) is totally real; \(K\) is imaginary quadratic.
Moreover \(H_0\subset\ker(\eta\eta')\), and
\(\eta\eta'\) is a nontrivial real quadratic character
(the two conductors are distinct).
Hence \([F_0:\mathbb Q]\) is even.

The CM extension \(L_0/F_0\) is unramified at every
finite place.
Away from the conductor of \(\eta\), its inertia is
already killed by \(\eta\).
At that conductor, \(\eta'\) is unramified, and the
condition \(\eta=\eta'\) on \(H_0\) kills the remaining
inertia.
Every prime in \(S\) splits in \(L_0/F_0\), since
\(\eta\) is trivial on its rational decomposition group.
At \(\ell\), \(\xi\) is unramified and the relation
\(\bar\Psi=\phi\) on \(H_0\) makes \(\bar\Psi\)
trivial on inertia, since the fresh \(\phi\) is
unramified there.
Thus (T2) is trivial on inertia over \(F_0\).
The extension \(F_0/\mathbb Q\) is tame at \(\ell\):
its inertia factors through the Teichm\"uller character,
whose order divides \(\ell-1\).

Multiplication by inertia at the three fresh conductor
primes changes \(\phi,\eta,\eta'\) independently while
leaving \(\xi,\Psi,\omega,\psi\) unchanged.
Indeed the latter characters are unramified at those
primes, and each of the other two auxiliary quadratic
characters is unramified at the prime used for a change.
This gives the explicit adjustment operation
\(\mathrm{adj2}\) in the source.

The element \(u\) of (T6) already lies in \(G_{L_0}\):
all three fresh quadratic characters are unramified
at \(2\), and \(\bar\Psi(u)=1\).
Apply the adjustment operation to complex conjugation,
setting \(\phi=1,\eta=\eta'=1\).
It gives an element \(v\in G_{L_0}\) with
\[
       \bar\Psi(v)=1,\qquad\bar\xi(v)=-1,\qquad
                                      \bar\omega(v)=-1 .
 \tag{T8}
\]
Since \(\psi\) is nontrivial of \(2\)-power order,
its image contains \(-1\).
The conductor of \(\psi\) is prime to \(\ell\);
Chinese remaindering in the cyclotomic quotient
therefore supplies an element with
\(\bar\psi=-1,\bar\omega=1\), and hence
\(\bar\Psi=-1\).
Adjust it to \(\phi=-1,\eta=\eta'=1\).
It lies in \(G_{L_0}\).
If its \(\bar\xi\)-value has square \(1\), multiply
it by \(u\).
Thus there is \(t\in G_{L_0}\) with
\[
       \bar\Psi(t)=-1,\qquad\bar\xi(t)^2\ne1,\qquad
                                      \bar\omega(t)=1 .
 \tag{T9}
\]

The witnesses \(u,v,t\) give all the character-separation
conditions used in the preceding subsections.
Every root ratio is
\(\bar\xi^x\bar\Psi^y\), with
\[
 (x,y)\in\{(\pm1,0),(0,\pm1),(\pm1,\pm1)\}.
\]
If \(x\ne0\), its value at \(u\) is nontrivial.
If \(x=0\), its value at \(t\) is \(-1\).
Hence every root character is nontrivial on \(G_{L_0}\);
this is \(\mathrm{Ok.g3}\).
These same witnesses belong to \(\ker\bar\omega\), so
the root characters remain nontrivial on that kernel.

More strongly,
\(\bar\omega\) is not in the subgroup generated by
all root ratios.
If
\(\bar\omega=\bar\xi^x\bar\Psi^y\) on \(G_{L_0}\),
evaluation at \(u\) makes \(x\) even, because
\(\bar\xi(u)\) has \(2\)-power order greater than four.
Evaluation at \(v\) then gives \(1=-1\), a contradiction.
This is \(\mathrm{Ok.g2}\) and supplies the premise of
Proposition~\ref{prop:sr-four-actual-scalar-place}.

The element \(t\) also has four distinct values in (T2):
they are
\(-\bar\xi(t),-1,\bar\xi(t),1\).
It can therefore be chosen as the frame's distinguished
regular element.
Finally, a character trivial on \(\ker\bar\omega\)
and permuting the four residual characters must be
trivial.
Compare its product with the character \(1\) in (T2);
it is a ratio of two of those characters.
Every nontrivial such ratio is nontrivial on
\(\ker\bar\omega\), as proved using \(u,t\).
This proves in particular the quadratic self-twist
exclusion \(\mathrm{Ok.g3'}\).

\subsubsection{The \(\ell\)-adic local witness and higher layers}
All three auxiliary quadratic characters split at \(\ell\).
The element \(g_\ell\) of (T5) therefore belongs to
\(G_{F_0}\), and to \(G_{L_0}\), while retaining its
rational local decomposition group.
Its values give
\[
          \bar\omega(g_\ell)=1,\qquad
          \bar\Psi(g_\ell)=1,\qquad
                                  \bar\xi(g_\ell)\ne1 .
 \tag{T10}
\]
These are exactly the second alternative of (D11),
hence the actual \(\mathrm{LocL}\) producer.
They are stronger than just a generic residual
non-equality of two characters.

For the larger dimension required by the support
calculation, adjoin successive finite layers of the real
cyclotomic \(\mathbb Z_\ell\)-extension.
Write \(F_m\) for the resulting real fields and
\(L_m=F_mK\).
The degrees \([F_m:\mathbb Q]\) grow without bound;
the extensions \(F_m/F_0\) have \(\ell\)-power degree.
Their finite unramified CM property and their splitting
at \(S\) follow by base change from \(L_0/F_0\).
The degree remains even.

In any such layer an element of \(G_{F_0}\) has an
odd \(\ell\)-power lying in \(G_{F_m}\).
Apply a common such power to \(u,v,t\).
Their \(\bar\xi,\bar\Psi\)-values have \(2\)-power order,
so all the displayed distinctions, signs and
root-separation arguments survive.
Thus \(\mathrm{g2}\), \(\mathrm{g3}\),
\(\mathrm{g3'}\), and the regular frame element survive
on every selected higher layer.
For the local witness, either take the same power of
\(g_\ell\), or note that its cyclotomic value is \(1\),
so it already fixes the cyclotomic layers.
Equation (T10) survives.
The source's odd-power transport and
\(\mathrm{leFit}\) implement exactly these facts.

This supplies the actual tower hypotheses used by
the scalar-place producer, selected-tower irreducibility
and the ordinary local nonlinking calculation.
All these witnesses come from the principal packet and
the explicitly chosen auxiliary dyadic character.
No Case III tower is used.

% Exact frozen-source correspondence:
% FrobeniusFields/ExistsDecompositionElementAtLXiPowFourNeOne,
% ExistsInertiaElementAtTwoXiPowFourNeOne, Adj2, Elts2,
% CharacterSubgroupOpenEvenIndexTamePower, Ok2B;
% CompactIntegralFields/ExistsTameTowerOkTwo;
% UpperTriangularPlaces/BigOm and NormalValuationBounds/NeOfPair.
% The ell-primary character Kap is trivial for this 2-power
% principal packet; the source permits it for larger packets.


\subsection{Producing the rank-four auxiliary scalar place}
\label{subsec:sr-four-scalar-place}
The auxiliary local vanishing condition used in the Selmer
calculation has an actual place producer.
It uses the full root-character condition of the selected
tower, rather than a guessed congruence for a rational prime.

\begin{srlemma}\label{lem:sr-finite-character-separation}
Let \(k\) be a finite field, \(\Gamma\) a group, and
\(\chi_0,\ldots,\chi_3,\omega:\Gamma\to k^\times\)
characters with
\(\chi_0\chi_3=\chi_1\chi_2\).
Suppose
\[
       \omega\ne(\chi_0/\chi_1)^u(\chi_0/\chi_2)^v
                            \quad\hbox{for every }u,v\in\mathbb Z .
 \tag{A1}
\]
Then there exists \(g\in\Gamma\) such that
\[
                     \chi_0(g)=\cdots=\chi_3(g),
                            \qquad\omega(g)\ne1.
 \tag{A2}
\]
\end{srlemma}
\begin{proof}
Replace \(\Gamma\) by its finite image under these five
characters, a finite abelian group \(A\) of exponent
dividing \(\#k^\times\).
Its character pairing with
\(A^\vee=\operatorname{Hom}(A,k^\times)\) is perfect.
Indeed the decomposition of \(A\) into cyclic factors
and the roots of unity in the cyclic group \(k^\times\)
prove this on each factor.
For a subgroup \(S\subset A^\vee\), the same decomposition
or the order formula for annihilators gives
\[
                       \operatorname{Ann}
                         (\operatorname{Ann}(S))=S.
\]
Let \(S\) be generated by all \(\chi_i/\chi_j\).
The paired relation expresses all of them as powers of
\(\chi_0/\chi_1\) and \(\chi_0/\chi_2\).
Hypothesis (A1) says \(\omega\notin S\).
Double annihilation therefore gives an element of
\(\operatorname{Ann}(S)\) not killed by \(\omega\).
It is exactly (A2).
\end{proof}

\begin{srproposition}\label{prop:sr-four-actual-scalar-place}
Let \(L/F\) be the CM extension of an actual Case II
tower, and restrict its residual diagonal characters
\(\bar\chi_i\) and cyclotomic character
\(\bar\epsilon_\ell\) to \(G_L\).
The tower condition
\[
 \bar\epsilon_\ell\ne
       (\bar\chi_0/\bar\chi_1)^u
       (\bar\chi_0/\bar\chi_2)^v
                   \quad\hbox{on }G_L
                         \quad(u,v\in\mathbb Z)
 \tag{A3}
\]
produces a finite place \(v_0\) of \(F\), split in \(L\),
outside any specified finite set, such that
the residual characters are unramified there and
\[
 \bar\chi_0|_{D_{v_0}}=\cdots=\bar\chi_3|_{D_{v_0}},
              \qquad
            \bar\epsilon_\ell|_{D_{v_0}}\ne1 .
 \tag{A4}
\]
Here the decomposition group is identified with that
of either chosen split place of \(L\).
Multiplying all four characters by the residual scalar
link character preserves (A4).
\end{srproposition}
\begin{proof}
The paired relation holds for the actual residual tuple.
Apply Lemma~\ref{lem:sr-finite-character-separation}
on \(G_L\) with \(\omega=\bar\epsilon_\ell\);
obtain \(g_0\) as in (A2).
All characters have open kernel.

Choose a finite Galois extension \(J/F\) containing \(L\)
on which these characters become trivial.
They are restrictions of the given rational characters,
so their common kernel is normal in \(G_F\).
The image of \(g_0\) in \(\operatorname{Gal}(J/F)\)
fixes \(L\).
Chebotarev supplies infinitely many places of \(F\),
unramified in \(J\), with this Frobenius conjugacy class.
Choose one outside the prescribed finite set.
It splits in \(L/F\), because its Frobenius fixes \(L\).
All the characters in question are unramified there.

Their values on a Frobenius lift are those at \(g_0\);
conjugation does not change a character.
Their values on the entire local decomposition group
are powers of these values, since unramified characters
factor through the procyclic unramified quotient.
The four values are equal on \(g_0\), hence on this
whole decomposition group.
The cyclotomic value at that Frobenius is nontrivial.
This proves (A4).
An arbitrary finite avoidance set allows the two links'
distinguished auxiliary places to be different.
A common scalar multiplication changes none of the
ratios or the cyclotomic assertion, proving the last claim.
\end{proof}

Condition (A3) is exactly \(\mathrm{Tower.Ok.g2}\), after
identifying the two root coordinates with the root weights.
The frozen source implements the finite-duality argument
in \srcid{WeakShimuraTheory/UnipotentIdempotents/GExRtZero.lean},
theorem \(\mathrm{Lk.gEx}\), and the actual place selection
in \srcid{EqOfKer.lean}, theorem \(\mathrm{Lk.setEx}\).
The version with an additional finite avoidance condition
is \(\mathrm{setAv}\) in
\srcid{Foundations/GaloisRepresentations/PlainTorsionPairs/LinkTwo.lean}.
The residual tuple in the Hida datum is the scalar link
twist of this tuple.
Thus \(\mathrm{Lk.v0ok}\) supplies exactly
\(\mathrm{r.V0ok}\), not merely a genericity statement
about residual characters.
The local vanishing transfer in
\(\mathrm{P1.n4V0}\) can therefore be applied with its
actual scalar-place premise.
No condition on a Frobenius at another local place
is substituted for (A4).


\subsection{Irreducibility on the actual CM tower}
\label{subsec:sr-actual-tower-irreducibility}
The irreducibility premise at the start of
Section~\ref{subsec:sr-case-two-descent} is produced by
the actual residual tuple.
The following argument also explains why no arbitrary
finite restriction is being assumed irreducible.

\begin{srlemma}\label{lem:sr-residual-distinct-abelian-restriction}
Let \(G\) be a group and \(H\) a normal subgroup of finite
index with \(G/H\) abelian.
Let \(V\) be an irreducible representation over an
algebraically closed field of characteristic zero,
with a stable lattice over a DVR.
Suppose the residual semisimplification is
\(\bigoplus_{i=1}^n\bar\chi_i\), where
the characters \(\bar\chi_i|_H\) are pairwise distinct.
Then \(V|_H\) is irreducible.
\end{srlemma}
\begin{proof}
The restriction \(V|_H\) is semisimple.
Choose a nonzero simple \(H\)-submodule of minimal
positive dimension.
Its \(G\)-translates are simple because \(H\) is normal;
their sum is all of \(V\) by irreducibility.
Successively adding these simple submodules gives a
semisimple module, proving the claim.

If \(V|_H\) were reducible,
\(\operatorname{End}_H(V)\) would have dimension greater
than one.
Conjugation by \(G\) acts on this algebra through the
finite abelian group \(G/H\).
In characteristic zero this action is a sum of characters,
and its invariant subspace is
\(\operatorname{End}_G(V)\), of dimension one by Schur.
There is therefore a nonzero eigenvector \(A\) for a
nontrivial character \(\eta\) of \(G/H\):
\[
                    \rho(g)A\rho(g)^{-1}=\eta(g)A .
\]
This is an intertwiner between two irreducible
\(G\)-representations, hence is invertible.
Choose \(g\) with \(\eta(g)\ne1\).
Since \(\eta|_H=1\), taking traces gives
\[
                      \operatorname{tr}\rho(gh)=0
                                     \quad(h\in H).
\]
Reduce this exact equality using the stable lattice.
The residual semisimplification has the same traces,
so
\[
              \sum_i\bar\chi_i(g)\bar\chi_i(h)=0
                                      \quad(h\in H).
 \tag{I1}
\]
Distinct multiplicative characters of a group are
linearly independent over their coefficient field.
For reference, in a shortest nonzero relation, translate
the argument by an element separating two characters
and subtract a scalar multiple of the original relation.
The remaining nonzero relation is shorter, a contradiction.
Apply this to \(\bar\chi_i|_H\) in (I1).
It gives \(\bar\chi_i(g)=0\) for every \(i\), impossible
for characters valued in a field's multiplicative group.
\end{proof}

\begin{srproposition}\label{prop:sr-case-two-tower-irreducible}
Every irreducible actual characteristic-zero rational
point of the Case II problem remains irreducible on
the CM field \(L\) of the selected tower.
\end{srproposition}
\begin{proof}
The compositum \(L=FK\) is abelian over \(\mathbb Q\):
\(F/\mathbb Q\) is abelian and \(K/\mathbb Q\) quadratic.
Thus \(H=G_L\) is normal of finite index with abelian
quotient.

The stable lattice in the actual point has residual
semisimplification
\(\bigoplus_{i=0}^3\bar\chi_i\).
The tower's root-separation condition is
\[
       \forall w\in\{\hbox{symplectic roots}\}\quad
       \exists g\in G_L:\quad \bar r_w(g)\ne1 .
 \tag{I2}
\]
For \(i\ne j\), the paired relation
\(\bar\chi_0\bar\chi_3=\bar\chi_1\bar\chi_2\)
expresses \(\bar\chi_i/\bar\chi_j\) as
\(\bar r_{w(i,j)}\) for a symplectic root.
Equation (I2) therefore says exactly that their
restrictions to \(G_L\) are pairwise distinct.
The possible relabeling of the ordinary tuple is a
permutation and changes none of these conclusions.
Apply Lemma~\ref{lem:sr-residual-distinct-abelian-restriction}.

If the supported representative is a conjugate of the
original lattice representation, first apply the lemma
to the original representation and then conjugate.
Irreducibility is unchanged.
Hence the exact premise of
Proposition~\ref{prop:sr-all-quadratic-automorphy}
has been supplied for every actual irreducible rational
point, and that proposition's conclusion gives
\(\mathrm{IsAutQA}\) for all of them.
\end{proof}

% Frozen source:
% DenseThresholdCases/MapRepComp.lean,
% private irr_sub, dist_L, irr_L, and theorem aut_n_alt.
% (I2) is exactly Tower.Ok.g3; the finite-abelian quotient
% is proved there by comm_L.


\subsection{Removing the link character and descending the actual point}
\label{subsec:sr-case-two-descent}
We finish the automorphic endpoint of the Case II support
argument. Keep the actual tower and Hecke quotients of
Section~\ref{subsec:sr-case-two-hecke}.
Let
\[
 r:G_{\mathbb Q}\longrightarrow
                   \mathrm{GSp}_4(\overline{\mathbb Q}_\ell)
 \quad\hbox{with multiplier }\theta
 \tag{D1}
\]
be the original regular arithmetic point.
Proposition~\ref{prop:sr-case-two-tower-irreducible} supplies
\(r|_{G_L}\) irreducible for the actual selected tower.
This is the precise hypothesis
in the source's \(\mathrm{pmAutO}\) and
\(\mathrm{pmDown}\).
It is not inferred by restricting an arbitrary irreducible
representation to an arbitrary finite extension.
The preceding support argument gives a regular algebraic
cuspidal representation attached to
\[
                         r|_{G_L}\otimes\kappa .
 \tag{D2}
\]
The two links, their common scalar character, and strong
multiplicity one give the asserted local factors away from
either distinguished auxiliary place, as proved there.

\begin{srlemma}\label{lem:sr-link-character-algebraic}
The actual link character \(\kappa\) is the \(\ell\)-adic
avatar of an algebraic Hecke character of \(L\).
Consequently (D2) can be untwisted to a regular algebraic
cuspidal representation attached exactly to \(r|_{G_L}\).
\end{srlemma}
\begin{proof}
Write the original twist as \(a=2q+b\), \(b\in\{0,1\}\).
The actual character construction gives
\[
 \kappa=\kappa_0\gamma^{-(a+4)}\beta^b,\qquad
 \gamma^2\omega=\epsilon_\ell,\qquad \omega^{\ell-1}=1,
 \qquad \kappa_0^m=1
 \tag{D3}
\]
for some positive \(m\).
At each place above \(\ell\) there are an integer \(j\)
and a positive integer \(k\) such that
\[
                 (\beta^2\epsilon_\ell^{-(2j+1)})^k=1
                          \quad\hbox{on inertia}.
\]
Put \(A_v=bj-q-2\), an integer. Direct multiplication
in this commutative character group gives
\[
              (\kappa\epsilon_\ell^{-A_v})
                          ^{\,2km(\ell-1)}=1
                           \quad\hbox{on }I_v .
 \tag{D4}
\]
For example, the cyclotomic exponent contributed by (D3)
is \(-(2q+b+4)/2+b(2j+1)/2=bj-q-2\);
multiplying by \(2km(\ell-1)\) eliminates all finite factors.
Equation (D4), rather than a choice of a half-cyclotomic
character, proves the assertion about integral local
weights.
The finite coefficient and continuity properties of the
actual character imply that the finite-order quotient in
(D4) has open kernel.
Thus \(\kappa\) is locally algebraic on units at every
place above \(\ell\), with integral, parallel exponent \(A_v\).
Its construction makes it unramified outside a finite set.

For completeness this local statement globalizes by
idele class reciprocity.
Let \(P_\ell\) be the algebraic product of the embeddings
of \(L\) into \(\overline{\mathbb Q}_\ell\) with these
integer exponents, using the sign determined by the chosen
reciprocity convention. Let \(P_\infty\) be the same product
of complex embeddings, transported by \(\iota\).
On finite ideles divide \(\kappa\circ\operatorname{Art}_L\)
by \(P_\ell\), and at infinity multiply by \(P_\infty\).
This defines
\[
 \chi(s_f,s_\infty)=
 \iota\!\left(
   \frac{\kappa(\operatorname{Art}_L(s_f))}
                    {P_\ell(s_\ell)}\right)
                         P_\infty(s_\infty).
 \tag{D5}
\]
The finite factor is locally constant: this is (D4) above
\(\ell\); away from \(\ell\) a continuous rank-one
\(\ell\)-adic character has finite image on local units
at all ramified places.
For \(\alpha\in L^\times\), global reciprocity kills
the principal finite idele after its connected complex
infinite component is included, and
\(P_\infty(\alpha)=\iota(P_\ell(\alpha))\).
Hence (D5) is trivial on \(L^\times\).
It is an algebraic Hecke character and its \(\ell\)-adic
avatar is precisely \(\kappa\).

Twisting the actual representation in (D2) by
\(\chi^{-1}\circ\det\) preserves cuspidality and regular
algebraicity and removes \(\kappa\) from its attached
Galois representation.
This gives the representation asserted in the lemma.
\end{proof}

We next spell out the elementary Galois correction in
solvable descent. It ensures exact attachment, including
the finite character which a bare descent statement may
leave unspecified.

\begin{srlemma}\label{lem:sr-cyclic-exact-descent}
Let \(M/K\) be a cyclic extension of prime degree between
CM fields, and let \(r_K:G_K\to
\mathrm{GL}_n(\overline{\mathbb Q}_\ell)\) be continuous,
with \(r_K|_{G_M}\) irreducible.
If \(r_K|_{G_M}\) is attached to a regular algebraic
cuspidal representation \(\Pi_M\), then \(r_K\) is attached
exactly to a regular algebraic cuspidal representation
\(\Pi_K\).
Conversely, regular algebraic cuspidal automorphy of
\(r_K\) implies such automorphy of its restriction, under
the same irreducibility hypothesis.
\end{srlemma}
\begin{proof}
For descent, \(\Pi_M\) is invariant under
\(\operatorname{Gal}(M/K)\).
Indeed its conjugate has the same unramified Frobenius
polynomials, since \(r_K\) extends the attached
representation; Chebotarev and strong multiplicity one
give the invariance.
The cyclic base-change and descent theorem of
Arthur--Clozel therefore gives a cuspidal descent
\(\Pi_K^0\).
At infinity every local extension is \(\mathbb C/\mathbb C\),
so its infinitesimal characters are the same algebraic
regular ones. Thus \(\Pi_K^0\) is regular algebraic.

To identify the Galois representation of this descent,
use
\href{https://virtualmath1.stanford.edu/~rltaylor/rigcoh.pdf}
{Harris--Lan--Taylor--Thorne, Theorem A}.
Its hypotheses are precisely a CM field, a regular
algebraic cuspidal general-linear representation and
an arbitrary coefficient prime.
In particular no polarization of the uncorrected
descent is needed.
It supplies the continuous semisimple representation
\(s=r_{\ell,\iota}(\Pi_K^0)\).
Unramified base change and Chebotarev give
\(s|_{G_M}\simeq r_K|_{G_M}\).

Conjugate \(s\) so that these restrictions agree.
For \(g\in G_K\), the operator \(s(g)r_K(g)^{-1}\)
commutes with their common irreducible \(G_M\)-action:
this follows by applying both extensions to \(ghg^{-1}\),
\(h\in G_M\).
Schur's lemma makes this operator a scalar \(\delta(g)\).
Multiplication shows that \(\delta\) is a character,
trivial on \(G_M\); it therefore has finite order and
factors through \(\operatorname{Gal}(M/K)\).
Use global class field theory to twist
\(\Pi_K^0\) by the inverse finite Hecke character.
The resulting \(\Pi_K\) is attached exactly to \(r_K\).
A finite twist preserves regular algebraicity and
cuspidality.

For ascent, cyclic base change of \(\Pi_K\) exists.
If it were not cuspidal, the cyclic cuspidality criterion
would give a nontrivial character \(\eta\) of
\(\operatorname{Gal}(M/K)\) with
\(\Pi_K\simeq\Pi_K\otimes\eta\).
Chebotarev would give \(r_K\simeq r_K\otimes\eta_\ell\).
An intertwiner for this self-twist commutes with \(G_M\),
hence is scalar by irreducibility on \(G_M\);
the intertwining identity would then force \(\eta=1\).
This contradiction proves cuspidality of the base change.
Its regular algebraicity is preserved at infinity,
and Chebotarev identifies the attached representation
with \(r_K|_{G_M}\).
\end{proof}

Only the following published base-change input is used here:
the cyclic prime-degree base-change/descent theorem and its
cuspidality criterion, in Arthur--Clozel,
\emph{Simple algebras, base change, and the advanced theory
of the trace formula}, Chapter 3.
Iterating the lemma along a subnormal series with
cyclic prime quotients proves the same exact descent
for a solvable Galois CM extension.
In our tower all intermediate fields containing \(K\)
are CM: the totally real field is abelian over
\(\mathbb Q\), and complex conjugation is central in
\(\operatorname{Gal}(L/\mathbb Q)\).

\begin{srproposition}\label{prop:sr-all-quadratic-automorphy}
The actual point (D1)--(D2) is regular algebraic cuspidal
automorphic over every imaginary quadratic field \(K'\)
for which \(r|_{G_{K'}}\) is irreducible.
\end{srproposition}
\begin{proof}
Apply Lemma~\ref{lem:sr-link-character-algebraic}
and solvable exact descent from \(L\) to its chosen
imaginary quadratic field \(K\).
The irreducibility on \(G_L\) implies irreducibility
on all larger Galois groups along this descent.
This gives an actual regular algebraic cuspidal
representation attached to \(r|_{G_K}\).

Fix \(K'\) with \(r|_{G_{K'}}\) irreducible.
By Lemma~\ref{lem:sr-quadratic-self-twists}, an irreducible
rank-four representation of \(G_K\) can become reducible
on at most \(16\) distinct quadratic extensions of \(K\).
As \(K''\ne K\) ranges over imaginary quadratic fields,
the extensions \(KK''/K\) are distinct.
Indeed a biquadratic field with two imaginary quadratic
subfields has a real third quadratic subfield.
Apply the same argument over \(K'\).
Among any fixed \(35\) distinct imaginary quadratic fields
there is therefore a field \(K''\), distinct from \(K,K'\),
for which
\[
                r|_{G_{KK''}},\qquad r|_{G_{K'K''}}
                         \quad\hbox{are irreducible}.
 \tag{D6}
\]
There are at most \(16+16+2=34\) excluded fields.
Cyclic ascent from \(K\) to \(KK''\), exact descent to
\(K''\), ascent to \(K'K''\), and exact descent to \(K'\)
now prove the assertion.
At every use of Lemma~\ref{lem:sr-cyclic-exact-descent},
the needed restriction is one of (D6).
In particular this argument never assumes that
\(r|_{G_{LK'}}\) is irreducible.
\end{proof}

\subsubsection{The actual polarization and full local comparison}
Let \(\Pi_{K'}\) be a representation obtained above.
The symplectic relation in (D1) and the fact that \(r\)
is defined on \(G_{\mathbb Q}\) give
\[
              (r|_{G_{K'}})^c
                    \simeq(r|_{G_{K'}})^\vee\theta .
\]
With the Galois normalization
\(\operatorname{rec}(\Pi|\det|^{(1-n)/2})\),
one has
\[
 r_{\ell,\iota}(\Pi^\vee)
       =r_{\ell,\iota}(\Pi)^\vee\epsilon_\ell^{1-n}.
\]
The polarizing Hecke character of \(\mathbb Q\) is therefore
the algebraic character with avatar
\(\theta\epsilon_\ell^{n-1}\), here
\(\theta\epsilon_\ell^3\).
The actual multiplier is locally algebraic and odd.
Thus its product with \(\epsilon_\ell^3\) has value \(1\)
at complex conjugation, the required rank-four parity.
Chebotarev and strong multiplicity one give the
corresponding automorphic polarization of \(\Pi_{K'}\).

We use full Frobenius-semisimplified local--global
compatibility for this regular algebraic essentially
conjugate-self-dual cuspidal representation.
This combined published scope is explicitly recorded in the
introduction to
\href{https://annals.math.princeton.edu/wp-content/uploads/annals-v179-n2-p03-p.pdf}
{Barnet--Lamb--Gee--Geraghty--Taylor, p.~507}.
Away from \(\ell\) it is the polarized realization theorem
there; at \(\ell\) its unrestricted monodromy comparison is
\href{https://arxiv.org/pdf/1202.4683}
{Caraiani, Theorem 1.1}, together with that polarized
normalization.
The latter theorem has no Iwahori or Shin-regular
hypothesis in its statement.
The comparison identifies \(N\), not only the
semisimplified Weil representation.
The attached \(r|_{G_{K'}}\) is de Rham above \(\ell\)
with the specified Hodge numbers.
These realization and comparison theorems apply to the
already automorphic point at every odd coefficient prime,
including rank four at \(\ell=3\);
the lower bounds on coefficient primes in automorphy
lifting theorems are not hypotheses here.

\subsubsection{The local window predicate}
We verify the complete local assertion called
\(\mathrm{WitA}\) in the source.
It has two clauses.
First, at \(v\mid\ell\), an ordinary flag with characters
\(\kappa_i\) satisfying
\[
                         \kappa_i\ne\kappa_j\epsilon_\ell
                                \quad\hbox{for all }i,j
 \tag{D7}
\]
can be replaced by a flag for which
\[
                    \kappa_i^m\ne
                      (\kappa_j\epsilon_\ell)^m
                       \quad(i,j,\ m\ge1).
 \tag{D8}
\]
These are \(\mathrm{Unl1}\) and
\(\mathrm{PotUnl}\), respectively.
Second, for every regular algebraic cusp representation
attached to \(r|_{G_{K'}}\), there is a single real number \(t\)
such that, at every \(v\nmid\ell\) and at every \(v\mid\ell\)
satisfying (D8), its cuspidal support consists of rank-one
characters \(\xi\) with
\[
                 q_v^{\,t-1/2}<|\xi(\varpi_v)|
                                  <q_v^{\,t+1/2}.
 \tag{D9}
\]
The source denotes this local assertion by
\(\mathrm{WinAt}(\Pi,t,v)\).
The normalizing scalar in Galois--automorphic comparison
shifts \(t\) and never changes a ratio.

\begin{srlemma}\label{lem:sr-parallel-wd-characters}
Let \(k/\mathbb Q_\ell\) be finite.
Among continuous de Rham characters of \(G_k\)
whose cyclotomic exponents are parallel integers,
the rank-one Weil--Deligne character determines the
Galois character.
Consequently, if a de Rham representation with such
rank-one ordinary graded characters satisfies (D7),
then its Weil--Deligne monodromy is zero.
\end{srlemma}
\begin{proof}
Write such a character as
\(\alpha=\epsilon_\ell^h\eta u\), with \(h\in\mathbb Z\),
\(\eta\) finite order and \(u\) unramified.
This follows from local reciprocity after dividing out
the parallel algebraic character on an open subgroup
of units.
The unramified value of \(u\) is an \(\ell\)-adic unit,
since \(\alpha\), \(\epsilon_\ell\), and \(\eta\)
are continuous characters of the compact group \(G_k\).

At a geometric Weil Frobenius, the valuation of
\(\mathrm{WD}(\alpha)\) is \(h\) times the nonzero
valuation of \(\mathrm{WD}(\epsilon_\ell)\).
The finite-order factor and the value of \(u\)
have valuation zero.
Thus the Weil--Deligne character first recovers \(h\).
After dividing out that power of \(\epsilon_\ell\),
it recovers the finite inertia and the unramified
value by local reciprocity.
This proves the first assertion.

The Weil--Deligne functor is exact on potentially
semistable representations.
For the ordinary flag its semisimplified Weil
representation therefore has the rank-one characters
\(\lambda_i=\mathrm{WD}(\kappa_i)\).
If \(N\ne0\), the Weil equivariance of monodromy
gives two of these characters with
\(\lambda_i=\lambda_j\mathrm{WD}(\epsilon_\ell)\),
after interchanging \(i,j\) if required by the
Frobenius convention.
The first assertion gives
\(\kappa_i=\kappa_j\epsilon_\ell\).
This contradicts (D7).
\end{proof}

The restriction to parallel exponents is essential to this
short proof and is satisfied here.
The rational ordinary point has cyclotomic exponents
\((a+1,a,1,0)\) on an open inertial subgroup.
Every restriction to a finite local field has the same
parallel exponents.
The ordinary flag over the real tower descends to
\(G_{\mathbb Q_\ell}\) as follows.
Pass to an open normal inertial subgroup where all finite
characters disappear.
Choose an element on which the cyclotomic character has
infinite order; it has the four distinct eigenvalues
\(\epsilon_\ell(g)^{a+1},
  \epsilon_\ell(g)^a,\epsilon_\ell(g),1\).
The ordered flag is the unique flag made from the
successive eigenspaces in this order.
Every local conjugate flag has the same ordered
characters on that subgroup, so equals it.
Hence the rational decomposition group preserves this
flag. Its graded characters have the stated parallel
exponents.
Restriction supplies such a flag at every place of every
imaginary quadratic field.

Away from \(\ell\), the actual point has finite abelian
inertia extending to decomposition characters.
Here is the precise ramification argument.
Use the fixed torus residual representation
\[
 \bar r_{\rm tor}
       =\operatorname{diag}(\bar\chi_0,\ldots,\bar\chi_3)
       =D.\mathrm{rbar},\qquad
                       H=G_F\cap\ker\bar r_{\rm tor}.
 \tag{D10}
\]
This is literally a diagonal representation.
In the original framed rational problem
\(\mathrm{KPair.red}\) gives entrywise reduction to
\(D.\mathrm{r0}=\operatorname{diag}(D.\mathrm{ps})\);
it is stronger than prescribing a semisimplification.
After the lattice changes used in the residual chain,
the paired ramification set \(\mathrm{fr.N}\) still uses
this same fixed \(D.\mathrm{rbar}\), rather than the
possibly nonsplit residual representation at a chain point.
Thus \(H\) is open and contains the commutator of
\(G_{\mathbb Q}\): \(F/\mathbb Q\) is abelian and
\(\bar r_{\rm tor}\) is diagonal in global characters.
The actual paired ramification condition kills
\(I_p\cap H\) for every \(p\ne\ell\).
Conjugation of the characteristic-zero point preserves
this killing assertion.
Thus \(r(I_p)\) is a finite abelian group, and conjugation
by a decomposition element is trivial on its quotient
through \(G_{\mathbb Q}/H\).
Each inertial isotypic subspace is decomposition-stable.
Diagonalizing a Frobenius on each such subspace shows
that the Frobenius-semisimplified Weil representation is
a sum of characters extending the inertial characters.
Finite inertia gives \(N=0\).
These statements remain true on every finite local
restriction.

\begin{srproposition}\label{prop:sr-actual-wit}
For every \(K'\) in Proposition~\ref{prop:sr-all-quadratic-automorphy},
the actual representation \(r|_{G_{K'}}\) satisfies
\(\mathrm{WitA}\).
Thus the endpoint is exactly the source predicate
\(\mathrm{IsAutQA}\), including its local assertion.
\end{srproposition}
\begin{proof}
A cuspidal automorphic representation of a general-linear
group becomes unitary after a single real norm twist.
The same twist works at all finite places.
Every local component of a cuspidal representation is
generic. The
\href{https://www.numdam.org/item/ASENS_1986_4_19_3_335_0/}
{generic unitary classification of Tadi\'c}
has the following consequence.
If its Weil--Deligne parameter is a sum of characters
and has \(N=0\), it is a principal series with inducing
characters in the strict window (D9).
Indeed its unitary building blocks are tempered
characters and complementary pairs with exponents
\(s,-s\), \(0<s<1/2\).
A discrete-series block of length greater than one
would have nonzero monodromy and is excluded.
Undoing the one global norm twist gives the one
number \(t\) in (D9).

At \(v\nmid\ell\), the finite abelian inertia calculation
after (D10) and full local--global compatibility give
exactly this character parameter with \(N=0\).
Hence (D9) holds there.

At \(v\mid\ell\), assume (D7).
There exists a flag with parallel graded exponents,
as just proved.
Its graded character multiset is the intrinsic multiset
of Jordan--H\"older characters of the local representation.
Any other full invariant flag gives the same multiset.
Thus (D7), if it holds for a flag, holds also for the
parallel flag.
The representation is de Rham by its actual polarized
automorphic realization, so
Lemma~\ref{lem:sr-parallel-wd-characters}
gives \(N=0\).
Its semisimplified Weil parameter is a sum of characters,
by exactness along the ordinary flag.
Full compatibility and the generic unitary classification
therefore give (D9).

Suppose now that (D8) fails.
Then for some \(i,j,m\ge1\),
\((\kappa_i/\kappa_j)^m=\epsilon_\ell^m\).
The Weil characters would have complex Frobenius
absolute-value ratio \(q_v\) or \(q_v^{-1}\);
the finite-order ambiguity has absolute value one.
Under local--global compatibility the common rank
normalizing scalar cancels, so the inducing characters
in (D9) would have ratio \(q_v\) in one order.
Their strict window makes that ratio strictly less
than \(q_v\), a contradiction.
The same parallel flag satisfies (D8).
This proves \(\mathrm{Unl1}\Rightarrow\mathrm{PotUnl}\).

Conversely (D8) implies (D7) by taking \(m=1\);
the preceding argument supplies (D9) at every such
place. This is exactly the second clause of
\(\mathrm{WitA}\).
Any other regular algebraic cusp representation
attached to this Galois representation equals
\(\Pi_{K'}\) by Chebotarev and strong multiplicity one,
so the asserted quantifier over all attached
representations introduces no extra hypothesis.
Together with
Proposition~\ref{prop:sr-all-quadratic-automorphy}
this proves \(\mathrm{IsAutQA}\).
\end{proof}

\subsubsection{The ordinary nonlinking condition at the selected tower}
We also record why the actual selected tower gives (D7)
before invoking the preceding implication.
This is the source's direct \(\mathrm{unlF}\) calculation;
it verifies a local hypothesis rather than assuming it
from generic residual characters.
The tower has the following alternative at every
\(\ell\)-adic place:
\[
 \begin{array}{ll}
 \text{either}&
 \exists g\in I_F:\ \bar\omega(g)\ne1,
                         \quad\bar r|_{I_F}=1,\\
 \text{or}&
 \exists g\in D_F:\quad
      \bar\omega(g)=1,\quad\bar\Psi(g)=1,\quad
                                     \bar\xi(g)\ne1 .
 \end{array}
 \tag{D11}
\]
Here \(\bar\omega=\bar\epsilon_\ell\),
\(\bar\Psi\) is the actual root character, and on
\(\ker\bar\Psi\) the residual diagonal is
\[
                      (\bar\xi,1,\bar\xi,1).
 \tag{D12}
\]
These are precisely the local tower condition
\(\mathrm{LocL}\) and the identities in
\(\mathrm{barFacts}\).

If \(\kappa_i=\kappa_j\epsilon_\ell\), restrict to an
open inertial subgroup where their finite factors
disappear.
The infinite cyclotomic image gives
\(h_i=h_j+1\).
For \((h_0,h_1,h_2,h_3)=(a+1,a,1,0)\), \(a\ge2\),
the only possibilities are
\((i,j)=(0,1),(2,3)\), or \((1,2)\), the last only
when \(a=2\).
Reduce the same character equality modulo the
coefficient maximal ideal.
In the first case of (D11) its two residual sides
are \(1\) and \(\bar\omega(g)\), a contradiction.
In the second case, (D12) has unequal \(i,j\)
entries for every listed pair, while
\(\bar\omega(g)=1\), again a contradiction.
Thus (D7) holds on the selected real tower.
The CM extension is split at these places, so the
same local decomposition groups give (D7) over \(L\).

Apply the local argument in
Proposition~\ref{prop:sr-actual-wit} to the actual
automorphic point over \(L\), undoing its algebraic
scalar link twist, whose ratios cancel.
It gives (D8) on \(G_{L_w}\).
An equality prohibited by (D8) on the rational
decomposition group would restrict to such an equality
on this finite-index subgroup.
Therefore (D8) holds already for the rational ordinary
flag and after every finite local restriction.
Indeed if the ratio became finite order on a
finite-index subgroup, a further positive power would
be trivial on the whole local group.
The actual point is consequently potentially crystalline
at \(\ell\), and has the required unlinked principal
series factors after each imaginary quadratic restriction.
This conclusion is obtained from actual automorphy,
the checked tower condition and full comparison.

The local--global comparison at a distinguished auxiliary
place uses the other link, as explained after (H16).
Thus no auxiliary prime or neat level is retained in this
endpoint.
This completes the actual Case II match from the common
ordinary module and its transported support to the full
automorphic-and-local predicate used by the counting
argument.

% Exact frozen-source counterparts:
% WeakShimuraTheory/UnipotentIdempotents/KapCyc;
% DenseThresholdCases/AuKap, TAlgOfCyc, TwistInvarianceAutomorphic;
% BaseChangeUnits/SolCOfSolDescent, VWOfOne, TAcross;
% DenseThresholdCases/PmAutO, DnPmDown;
% UpperTriangularPlaces/PtLawsOf and
% NormalValuationBounds/NeOfPair, CommLeRbar;
% AutomorphyLifting/WeakOfAtt and IsAutCsdA.
% The selected-tower irreducibility premise is explicit in (D1).
